\documentclass{situationist-preprint}

\title{Existence Precedes Essence}
\msid{EXI-09-a}
\msfamily{EXI-09}
\msdiscipline{Meaning, Death and the Absurd}
\mstype{Principal manuscript}
\mskeywords{existentialism; Sartre; essence; artefacts; paper-knife; maker's intention; project; real definition}

\begin{document}
\maketitle

\begin{abstract}
Sartre held that for human beings existence precedes essence, while for an artefact such as a paper-knife essence precedes existence: the artisan conceives the knife and what it is for before making it. The claim is often dismissed, either as resting on atheism or as contradicting the thesis that every object, human beings included, has an essence. We argue that Sartre's distinction holds. The key move is to read Sartre's \enquote{essence} as a defining project, a determination of what a thing is for that organizes its existence, and not as a real definition stating what kind of thing it is. On this reading, an artefact's essence is fixed by its maker's intention prior to the artefact's existence, because an artefact's kind is fixed by that intention; and a human being's essence is fixed by no intention prior to her existence, because a defining project for a life can be fixed only by the agent who lives it. The result does not depend on atheism, and it is compatible with the thesis of MET-13 that human beings, like everything else, have real definitions.
\end{abstract}

\section{The problem}

In \emph{Existentialism Is a Humanism} Sartre introduces his slogan by contrast with an artefact \citep{Sartre1946}. A paper-knife is made by an artisan who has a conception of it, of the technique for producing it and of the use it is to serve. Its essence, the set of procedures and properties that define it and allow it to be made, precedes its existence. If God created human beings, they would be like paper-knives, made according to a conception in the divine understanding. Sartre's atheist existentialism denies this: there is at least one being whose existence comes before its essence, a being that exists before it can be defined by any concept. Man first exists, encounters himself, and defines himself afterwards; he is nothing other than what he makes of himself.

Two objections are standard. The first is that the slogan rests on atheism, so that a theist has an immediate reply. The second is that it contradicts essentialism: if every object has a real definition, as MET-13 argues, human beings have essences that do not wait upon their choices. A third, less often stated, is that the artefact side of the contrast is too quick, since artefacts acquire functions their makers never intended. We argue that all three objections misidentify what Sartre means by essence, and that once this is fixed his distinction holds: for human beings existence precedes essence, and for artefacts it does not.

\section{Framework}

\begin{definition}[Defining project]\label{d:proj}
A \emph{defining project} for an entity $x$ is a determination of what $x$ is for, in the light of which the activities, uses or changes of $x$ count as succeeding or failing as what $x$ is.
\end{definition}

\begin{definition}[Precedence]\label{d:prec}
For an entity $x$, \emph{essence precedes existence} if a defining project for $x$ is fixed before $x$ exists, by a determination on which $x$'s coming into existence depends. \emph{Existence precedes essence} if no defining project for $x$ is fixed before $x$ exists, and whatever defining project $x$ has is fixed in the course of $x$'s existence.
\end{definition}

Definition~\ref{d:proj} is our reading of Sartre's term. When he says that the paper-knife's essence precedes it, he lists its production method and its use; when he says that man defines himself afterwards, he means that a person's life comes to be for something only through what she undertakes \citep{Sartre1946}. A defining project is not a real definition in the sense of MET-13, which states what kind of thing an object is and its origin. The two can coincide in their content, as for artefacts, but they need not, as for human beings.

\begin{principle}[Artefact kinds]\label{p:art}
An artefact belongs to its kind in virtue of having been made by someone who intended to make a thing of that kind, with a substantive conception of what things of that kind are for, and who largely succeeded \citep{Hilpinen1992,Thomasson2003}.
\end{principle}

\begin{principle}[Projects of agents]\label{p:agent}
For a being capable of adopting ends, what its life is for, in the sense of Definition~\ref{d:proj}, is fixed by the ends it adopts and pursues. Another's intention for that life can give it reasons, but does not fix its defining project.
\end{principle}

Principle~\ref{p:art} is the standard intentional account of artefact kinds, and MET-13 adopts it for artefacts: the real definition of an artefact states its kind, which is fixed by the function its maker intended. Principle~\ref{p:agent} states that the success conditions of a life, as opposed to those of a tool, are internal to the agent's own ends: a person whose parents intended her to be a physician and who becomes a composer has not failed as what she is, whereas a paper-knife that cannot cut paper has failed as a paper-knife. The principle is the core of Sartre's claim that man is what he makes of himself, and it is close to Korsgaard's account of practical identity as a self-conception under which one values oneself and finds one's life worth living \citep{Korsgaard1996}.

\section{The argument}

\begin{lemma}\label{l:art}
For every artefact, essence precedes existence.
\end{lemma}

\begin{proof}
Let $x$ be an artefact of kind $K$. By Principle~\ref{p:art}, $x$ is of kind $K$ because a maker intended to make a $K$ with a conception of what $K$s are for. That conception is a determination of what $x$ is for, in the light of which $x$'s uses count as succeeding or failing (a knife that does not cut is a defective knife), so it is a defining project for $x$ (Definition~\ref{d:proj}). The intention was formed before $x$ existed, since $x$ was made in carrying it out, and $x$'s coming into existence as a $K$ depends on it. By Definition~\ref{d:prec}, essence precedes existence for $x$.
\end{proof}

\begin{lemma}\label{l:hum}
For every human being, existence precedes essence.
\end{lemma}

\begin{proof}
Let $h$ be a human being. Before $h$ exists, there may be intentions concerning $h$: those of parents, of a society, or, on a theistic hypothesis, of a creator. By Principle~\ref{p:agent}, none of them fixes a defining project for $h$, since $h$ is capable of adopting ends and what $h$'s life is for is fixed by the ends $h$ adopts. Nor does $h$'s kind fix one: the kind \emph{human organism capable of thought} is fixed by biological origin, not by an intention, and it is compatible with indefinitely many incompatible lives, none of which fails as a human life by departing from the others. So no defining project for $h$ is fixed before $h$ exists. Whatever defining project $h$ has is fixed by the ends $h$ adopts, which $h$ can adopt only in the course of existing. By Definition~\ref{d:prec}, existence precedes essence for $h$.
\end{proof}

\begin{theorem}\label{t:main}
Sartre's distinction holds: for human beings existence precedes essence, and for artefacts essence precedes existence.
\end{theorem}

\begin{proof}
Immediate from Lemmas~\ref{l:art} and~\ref{l:hum}.
\end{proof}

\begin{corollary}\label{c:met}
Theorem~\ref{t:main} is compatible with the thesis that every object, including every human being, has a real definition.
\end{corollary}

\begin{proof}
For artefacts, the real definition states the kind fixed by the maker's intention, which by Lemma~\ref{l:art} is also prior to the artefact's existence, so the two notions of essence agree in content and in order. For human beings, the real definition states the kind and the origin, and is fixed when the human being comes to exist; Lemma~\ref{l:hum} concerns defining projects, which the real definition does not supply. Nothing in the real definition of a human being settles what her life is for.
\end{proof}

Corollary~\ref{c:met} records the position MET-13-a takes on the relation between its essentialism and the present result: Sartre's essence is a defining project, and his claim is that no such project is given to human beings in advance of their existence.

\section{Objections and replies}

\begin{objection}[Theism]\label{o:god}
Sartre derives his thesis from atheism, as he says himself. If God made human beings according to a conception, they are like paper-knives.
\end{objection}

\begin{reply}
Lemma~\ref{l:hum} does not use atheism. It grants that intentions concerning $h$ may precede $h$, including a creator's, and denies only that such intentions fix a defining project for a being capable of adopting ends. A creator who makes rational agents makes beings whose lives are for what they take up; a creator's purpose can give them weighty reasons, but a person who rejects it has not thereby failed to be a person, as a knife that does not cut has failed to be a knife. Sartre's atheism gives his slogan its rhetorical force; Principle~\ref{p:agent} gives it its truth.
\end{reply}

\begin{objection}[Human nature]\label{o:arist}
Aristotle argued that the human being has a characteristic function, activity of soul in accordance with reason, and that the human good is fixed by it \citep[I.7]{AristotleNE}. Human beings therefore have a defining project prior to any choice.
\end{objection}

\begin{reply}
If the function argument succeeds, it fixes a general form that any good human life takes, not what a particular life is for. Rational activity is realized equally by the lives of the composer and the physician, and the argument gives no ground for saying that either has failed as what she is. A defining project in the sense of Definition~\ref{d:proj} is the more specific determination that the general form leaves open, and it is that determination which, on Sartre's view and ours, waits upon existence.
\end{reply}

\begin{objection}[Social formation]\label{o:soc}
Human beings are born into roles and expectations that shape what they become, often decisively. \citet{Beauvoir1949} shows how a social destiny is imposed on women before they can choose.
\end{objection}

\begin{reply}
Beauvoir's analysis supports the thesis rather than threatening it. Her point is that the destiny presented as feminine nature is a social imposition and not a defining project that a woman fails by refusing. Social roles that precede a person are possibilities and pressures; by Principle~\ref{p:agent}, they become her defining project only if she adopts them.
\end{reply}

\begin{objection}[Repurposing]\label{o:rep}
Artefacts acquire functions their makers did not intend. A piece of driftwood becomes a table when someone puts a plate on it, and a paper-knife used only as a bookmark is for marking pages. The artefact's essence therefore does not precede its existence.
\end{objection}

\begin{reply}
In the driftwood case, what existed before the intention was a piece of wood, not a table; the table comes into existence, as an artefact of its kind, through the intentional act of selecting and placing the wood, and that intention precedes it, as Principle~\ref{p:art} requires \citep{Hilpinen1992}. In the bookmark case, the paper-knife acquires a use, not a new kind: it remains a paper-knife used as a bookmark, and it still fails as a paper-knife if it cannot cut. Where a new artefact kind is created by a new intention, the new artefact's essence precedes its existence as that artefact.
\end{reply}

\section{What the result does not establish}

The result concerns artefacts and beings capable of adopting ends. It is silent on organisms and natural objects that are neither, which have real definitions but no defining projects in our sense. It does not establish Sartre's further theses that freedom is absolute, that each choice legislates for humanity, or that there are no values prior to choice. Principle~\ref{p:agent} says that what a life is for is fixed by its agent's ends; it does not say that every end is as good as any other, and is consistent with EXI-01, on which a life is meaningful to the extent that its projects are of objective worth.

\section{Conclusion}

Sartre's distinction holds. Read as a defining project, an artefact's essence is fixed by the intention that brings it into being and so precedes it, while a human being's defining project can be fixed only by the ends she adopts and so comes after her existence. The argument does not depend on atheism, and it leaves intact the thesis that human beings, like everything else, have real definitions, which state what they are without settling what their lives are for.

\begin{thebibliography}{9}\small
\bibitem[Aristotle(NE)]{AristotleNE} Aristotle. \emph{Nicomachean Ethics}. Cited by book and chapter.
\bibitem[Beauvoir(1949)]{Beauvoir1949} Beauvoir, S. de (1949). \emph{Le Deuxi\`eme Sexe}. Gallimard.
\bibitem[Hilpinen(1992)]{Hilpinen1992} Hilpinen, R. (1992). Artifacts and works of art. \emph{Theoria}, 58(1), 58--82.
\bibitem[Korsgaard(1996)]{Korsgaard1996} Korsgaard, C.~M. (1996). \emph{The Sources of Normativity}. Cambridge University Press.
\bibitem[Sartre(1946)]{Sartre1946} Sartre, J.-P. (1946). \emph{L'existentialisme est un humanisme}. Nagel.
\bibitem[Thomasson(2003)]{Thomasson2003} Thomasson, A.~L. (2003). Realism and human kinds. \emph{Philosophy and Phenomenological Research}, 67(3), 580--609.
\end{thebibliography}

\end{document}
